\subsection{Similarity invariants of an endomorphism}

If we translate the decomposition of a finitely generated torsion module in VII, p.~8, Th. 1 and p.~9, Prop. 2, then we obtain :

PROPOSITION 3. --- Let E be a vector space of finite dimension n over a commutative field K, and let u be an endomorphism of E; for every monic irreducible polynomial \(p(\mathbf{X})\) , let \(M_{p}\) be the vector subspace consisting of elements x of E such that \((p(u))^{k}\) . x = 0 for some integer k. Then \(M_{p}\) is closed under u, the vector space E is the direct sum of the \(M_{p}\) , and there exist polynomials \(s_{p}\) such that, for all \(x \in E\) , the component of x in \(M_{p}\) is equal to \(s_{p}(u)\) . x.

Remark 1. --- Clearly the minimal polynomial of the restriction of u to \(M_{p}\) is the greatest power of p which divides the minimal polynomial of u. Moreover we have \(s_{p}(u).x = x\) for \(x \in M_{p}\) , from which it follows immediately that, if \(M_{p} \neq 0\) , then \(s_{p}\) is coprime to p.

Similarly, by Th. 2 of VII, p.~19, the module \(\mathbf{E}_u\) is isomorphic to a direct sum of cyclic modules \(\mathrm{F}_j = \mathrm{K}[\mathrm{X}] / \mathfrak{a}_j\) ( \(1 \leqslant j \leqslant r\) ), where the ideals \(\mathfrak{a}_j\) are distinct from \(\mathrm{K}[\mathrm{X}]\) and \(\mathfrak{a}_j \subset \mathfrak{a}_{j+1}\) ; and the \(\mathfrak{a}_j\) are determined by these conditions. Moreover, since \(\mathrm{E}_u\) is a torsion module, we have \(\mathfrak{a}_1 \neq (0)\) ; since E has dimension n, we have \(r \leqslant n\) . Put \(\mathfrak{a}_j = (h_j)\) ( \(1 \leqslant j \leqslant r\) ), with \(h_j\) a monic polynomial, and consider the sequence \((q_i)\) ( \(1 \leqslant i \leqslant n\) ) of polynomials defined by:

\[
\left\{ \begin{array}{l l} q _ {i} (\mathbf {X}) = 1 & \text { if } \quad i \leqslant n - r \\ q _ {i} (\mathbf {X}) = h _ {n - i + 1} (\mathbf {X}) & \text { if } \quad n - r <   i \leqslant n. \end{array} \right.\tag{3}
\]

It is clear that the polynomials \(q_{i}\) determine the polynomials \(h_{j}\) and conversely, and that \(E_{u}\) is isomorphic to the direct sum of the n modules \(\mathrm{K}[\mathrm{X}]/(q_{i})\) , the first n-r of which are 0.

In other words :

PROPOSITION 4. --- Let E be a vector space of finite dimension n over a commutative field K, and let u be an endomorphism of E. Then there exist n monic polynomials \(q_{i}(X) \in K[X]\) ( \(1 \leqslant i \leqslant n\) ) such that \(q_{i}\) divides \(q_{i+1}\) for \(1 \leqslant i \leqslant n-1\) , and E is the direct sum of n subspaces \(V_{i}\) ( \(1 \leqslant i \leqslant n\) ) closed under u, cyclic (with respect to u), and such that the minimal polynomial of the restriction of u to \(V_{i}\) is equal to \(q_{i}\) ( \(1 \leqslant i \leqslant n\) ). The polynomials \(q_{i}\) are uniquely determined by these conditions, and \(q_{n}\) is the minimal polynomial q of u.

Remark 2. --- By the above proposition, there exists a basis of E with respect to which the matrix U of u has the form

\[
\left( \begin{array}{c c c c c} A _ {n - r + 1} & 0 & \dots & 0 & 0 \\ 0 & A _ {n - r + 2} & \dots & 0 & 0 \\ \cdots & \cdots & \cdots & \cdots & \cdots \\ 0 & 0 & \dots & A _ {n - 1} & 0 \\ 0 & 0 & & 0 & A _ {n} \end{array} \right)
\]

where each matrix \(A_{i}\) has the form (2) (taking \(g(\mathbf{X}) = q_{i}(\mathbf{X})\) ) (cf.~VII, pp.~29-30).

DEFINITION 2. --- In the notation of Prop. 4, the \(n\) monic polynomials \(q_{i}(\mathbf{X}) (1 \leqslant i \leqslant n)\) are called the similarity invariants of the endomorphism \(u\) .

Thus the n-th similarity invariant \(q_{n}\) is the minimal polynomial of u (Prop. 4); in other words, for a polynomial \(p(\mathbf{X}) \in \mathbf{K}[\mathbf{X}]\) to satisfy \(p(u) = 0\) , it is necessary and sufficient that p be a multiple of \(q_{n}\) .

COROLLARY 1. --- Let K be a field, let E and E' be two finite dimensional vector spaces over K, and let u (resp. u') be an endomorphism of E (resp. E'). Then u and u' are similar (VII, p.~29) if and only if they have the same similarity invariants. Indeed u and u' are similar if and only if the K{[}X{]}-modules \(E_{u}\) and \(E_{u'}\) are isomorphic.

COROLLARY 2. --- Let u be an endomorphism of a finite dimensional vector space E over a field K, let \((q_{1}, \ldots, q_{n})\) be the family of similarity invariants of u, let L be an extension of K, let \(\mathrm{E}_{(\mathrm{L})} = \mathrm{L} \otimes_{\mathrm{K}} \mathrm{E}\) be the L-vector space induced from E by extension of scalars and let \(u_{(\mathrm{L})} = 1_{\mathrm{L}} \otimes u\) be the endomorphism of \(\mathrm{E}_{(\mathrm{L})}\) induced by u. Then the similarity invariants of \(u_{(\mathrm{L})}\) are the images \(\bar{q}_{1}, \ldots, \bar{q}_{n}\) of \(q_{1}, \ldots, q_{n}\) in L{[}X{]}.

This follows immediately from Prop. 4 and the fact that the L{[}X{]}-modules \(\mathrm{E}_{(\mathrm{L})u_{(\mathrm{L})}}\) and \((\mathrm{K}[\mathrm{X}]/(q_{i}))_{(\mathrm{L})}\) are isomorphic to L{[}X{]}⊗ \(_{K[X]}\) \(E_{u}\) and L{[}X{]}/( \(\bar{q}_{i}\) ) respectively.

Let U be a square matrix of order n with coefficients in a commutative field K. Then the similarity invariants of the endomorphism of \(K^{n}\) defined by U are called the similarity invariants of U. It then follows from Cor. 1 above that two square matrices are similar if and only if they have the same similarity invariants, and that if u is an endomorphism of a finite dimensional vector space E over K, and U is the matrix of u with respect to some basis B of E, then the similarity invariants of U and u coincide. By Cor. 1 and 2 above, we have :

COROLLARY 3. --- Let U and V be two square matrices of order n with coefficients in a commutative field K. If there exists an invertible square matrix P over some extension \(K'\) of K such that \(V = P^{-1}UP\) , then there exists an invertible square matrix Q over K such that \(V = Q^{-1}UQ\) .

Let E be a finite dimensional vector space over a commutative field K, let \((e_{i})_{1 \leq i \leq n}\) be a basis of E and let u be an endomorphism of E. Then by the exact sequence (1) of VII, p.~29, the K{[}X{]}-module \(E_{u}\) associated to u is isomorphic to the quotient of the free K{[}X{]}-module E{[}X{]}, with basis \((1 \otimes e_{i})\) , by the image of E{[}X{]} under the K{[}X{]}-linear map X - \(\bar{u}\) . The similarity invariants \(q_{i}(X)\) of u (VII, p.~32, Def. 2) are thus the invariant factors of X - \(\bar{u}\) (VII, p.~22). Thus Prop. 6 of VII, p.~22, implies :

PROPOSITION 5. --- Let E be a vector space of finite dimension n over a commutative field K, let u be an endomorphism of E, and let U be its matrix with respect to some basis of E. Then for each integer m with \(1 \leq m \leq n\) , the product

\[
d _ {m} (\mathbf {X}) = q _ {1} (\mathbf {X}) q _ {2} (\mathbf {X}) \dots q _ {m} (\mathbf {X})
\]

of the first \(m\) similarity invariants of \(u\) is equal to the gcd of the \(m\) -th order minors of the matrix \(XI_{n} - U\) .

COROLLARY 1. --- Let u be an endomorphism of a vector space of finite dimension n over a commutative field K, with characteristic polynomial \(\chi_{u}(\mathbf{X})\) and similarity invariants \(q_{i}(\mathbf{X})\) ( \(1 \leqslant i \leqslant n\) ). Then

\[
\chi_ {u} (\mathbf {X}) = q _ {1} (\mathbf {X}) q _ {2} (\mathbf {X}) \dots q _ {n} (\mathbf {X}).
\]

COROLLARY 2. --- In the notation of Cor. 1, let \(q(\mathbf{X})\) be the minimal polynomial of \(u\) ; then \(q(\mathbf{X})\) divides \(\chi_u(\mathbf{X})\) and \(\chi_u(\mathbf{X})\) divides \(q(\mathbf{X})^n\) . In particular the minimal polynomial and the characteristic polynomial of \(u\) have the same roots, and these are the eigenvalues of \(u\) .

Since \(q(\mathbf{X}) = q_{n}(\mathbf{X})\) , it is clear that \(q(\mathbf{X})\) divides \(\chi_{u}(\mathbf{X})\) . On the other hand, since each \(q_{i}\) divides q, their product \(\chi_{u}\) divides \(q^{n}\) .

COROLLARY 3. --- An endomorphism \(u\) is nilpotent if and only if its characteristic polynomial has the form \(X^n\) .

This follows immediately from Cor. 2.

Let us now rewrite Prop. 9 of VII, p.~24, which gives the decomposition of a module as a direct sum of indecomposable submodules.

PROPOSITION 6. --- Let E be a vector space of finite dimension n over a commutative field K, and let u be an endomorphism of E. Then E is the direct sum of subspaces \(E_{k}\) , closed under u and cyclic with respect to u, such that the minimal polynomial of the restriction of u to \(E_{k}\) has the form \(p_{k}^{n(k)}\) , where \(p_{k}\) is an irreducible polynomial, and \(E_{k}\) cannot be expressed as a direct sum of two nonzero subspaces each closed under u. For every monic irreducible polynomial \(p \in K[X]\) and every integer \(n \geqslant 1\) , the number \(m(p^{n})\) of subspaces \(E_{k}\) in any such decomposition, such that \(p^{n}\) is the minimal polynomial of the restriction of u to \(E_{k}\) , is uniquely determined.

The \(p_{k}^{n(k)}\) determine the similarity invariants of u and vice versa; we can get from one to the other by the procedure explained in VII, p.~25, Remarks 2 and 3. Furthermore, we can immediately get from the decomposition considered in Prop. 6 to those considered in Prop. 3 and 4.

Note that the monic irreducible polynomials \(p \in K[X]\) such that \(m(p^{n}) > 0\) for some integer \(n \geqslant 1\) are precisely the monic irreducible factors of the minimal polynomial of u. Thus, in contrast to the similarity invariants, these polynomials depend in general on the field K in which we are working.

